\section{Background on measure theory}

Let $X$ be a set.

\begin{definition}[Semi-ring, ring, $\sigma$-ring]
  A family $\mathcal{S}$ of subsets of a set $X$ is called a
  \emph{semiring} (of sets) if it contains the empty set,
  $A \cap B \in \mathcal{S}$ for all $A,B \in \mathcal{S}$ and, for
  every pair of sets $A,B \in \mathcal{S}$ with $A \subset B$, the set
  $B \backslash A$ is a disjoint union of finitely many sets
  in~$\mathcal{S}$. A family $\mathcal{R}$ of subsets of a set $X$ is
  called a \emph{ring} if it contains the empty set and for all
  $A,B \in \mathcal{R}$, the sets $A \cap B$, $A \cup B$ and
  $A \backslash B$ belong to~$\mathcal{R}$. If $\mathcal{R}$ is a
  ring, the \emph{$\sigma$-ring} of $\mathcal{R}$ is the set of
  countable unions of elements of~$\mathcal{R}$.
\end{definition}

\begin{definition}[semialgebra, algebra, $\sigma$-algebra]
  If in the above definitions one assumes $X$ is also an element of
  the family, one gets, respectively, the definitions of
  \emph{semialgebra}, \emph{algebra} and \emph{$\sigma$-algebra}.
\end{definition}

\begin{lemma}
  For any semiring $\mathcal{S}$, the collections of all finite unions
  of sets in $\mathcal{S}$ forms a ring $\mathcal{R}_{\mathcal{S}}$,
  the \emph{ring generated by $\mathcal{S}$}. Every set in
  $\mathcal{R}_{\mathcal{S}}$ is a finite union of pairwise disjoint
  sets in~$\mathcal{S}$.
\end{lemma}
\begin{proof}
See \cite[1.2.14~Lemma]{Bo07:MeasureTheory}.
\end{proof}

\begin{definition}
A real-valued set function $\mu$ defined on a class of sets $\mathcal{A}$ is called \emph{finitely additive} if 
\[
\mu\left(\bigcup_{i=1}^n A_i \right) = \sum_{i=1}^n \mu(A_i)
\]
for all $n$ and all disjoint sets $A_1,\cdots,A_n \in \mathcal{A}$ such that $\bigcup_{i=1}^n A_i \in \mathcal{A}$.
$\mu$ is called \emph{countably subadditive} if 
\[
\mu\left(\bigcup_{i=1}^{\infty} A_i \right) = \sum_{i=1}^{\infty} \mu(A_i)
\]
for all pairwise disjoint sets $A_i \in \mathcal{A}$ such that $\bigcup_{i=1}^{\infty} A_i \in \mathcal{A}$.
We say $\mu$ is called \emph{countably additive} if the above inequality is replaced by an equality.
\end{definition}

\begin{proposition}
A finitely additive function $\mu$ on a semiring $\mathcal{S}$ extends
uniquely to a finitely additive function on $\mathcal{R}_{S}$.
\label{prop:addsemitoring}
\end{proposition}
\begin{proof}
See \cite[1.3.10~Proposition]{Bo07:MeasureTheory}.
\end{proof}


\begin{definition}
  Fix a ring $\mathcal{R}$. Let $\mathcal{F}$ be a family of nested sequences
  of elements of $\mathcal{R}$. Then $\mathcal{F}$ is a \emph{final family} if,
  for any sequence $(R_i)_{i=0}^{\infty}$ of nested elements of
  $\mathcal{R}$ with $\bigcap_{i=0}^\infty R_i = \emptyset$, there is a sequence
  $(S_j) \in \mathcal{F}$ so that $R_i$ is eventually contained
  in~$S_j$, in the sense that for each $j$ there is an $N$ so that for
  all $i \ge N$ we have $R_i \subset S_j$.
  \label{def:finalfamily}
\end{definition}

To show countable additivity, it suffices to show that there is a final
family on which measures go to zero.


\begin{proposition}
Let $\mu$ be a finitely additive function defined on a ring $\mathcal{R}$. Let $\mathcal{F}$ be a final family of $\mathcal{R}$ so that for any $(S_j)_{j=1}^{\infty} \in \mathcal{F}$, $\lim_j \mu(S_j) \to 0$. Then $\mu$ is countably subadditive.
\label{prop:finalfamcount}
\end{proposition}
\begin{proof}
Let $A_n$ be a sequence of pair-wise disjoint sets in $\mathcal{R}$ whose union $A$ is also in $\mathcal{R}$.
Then consider $B_n =  A - \cup_{i=1}^n A_i$ . The sequence $B_n$ is a nested decreasing sequence of elements in $\mathcal{R}$ with empty intersection. Since $\mathcal{F}$ is a final family, for every $j$ we can find $N$ so that for every $i \geq N$, we have $B_i \subset S_j$.
On the one hand, note that $\mu(B_i) \to 0$, since $\mu(S_j) \to 0$ by definition of final family and $\mu$ being monotone.
On the other hand, 
\[
\mu(A - \cup_{i=1}^n A_i )= \mu(A) - \mu( \cup_{i=1}^n A_i )
\]
by finite additivity of $\mu$, which proves what we wanted.
\end{proof}


\begin{definition}
\label{def:outermeasure}
A set function $\mu$ defined on the class of all subsets of a set $X$ and taking values in $[0,\infty]$ is called an \emph{outer measure} if

\begin{enumerate}
    \item $\mu(\emptyset)=0$
    \item $\mu(A) \leq \mu(B)$ whenever $A \subset B$, i.e., is monotone
    \item $\mu$ is countably subadditive. 
\end{enumerate}
A set $A \subset X$ is called $\mu$-measurable if, for every $E \subset X$, one has the equality
\[
\mu(E \cap A) + \mu(E \backslash A) = \mu(E)
\]
The class of all $\mu$-measurable sets is denoted by $\mathcal{M}_{\mu}$, and is an algebra \cite[Theorem~1.11.4]{Bo07:MeasureTheory}.
An outer measure is called \emph{regular} if, for every set $A \subset X$, there exist a $\mu$-measurable set $B$ so that $A \subset B$ and $\mu(A)=\mu(B)$.
\end{definition}

\begin{theorem}[Caratheodory extension theorem]
Given a ring $\mathcal{R}$ and a finitely additive and countably additive function $\mu$ on it.
Let 
\[
\mu^*(A) \coloneqq \inf \{ \sum_{i=1}^n \mu(A_n) \bigm\vert A_n \in \mathcal{R}, A \subset \cup_{i=1}^{\infty} X_i\},
\]
where in the case of absence of such sets $A_n$ we set $\mu^*(A) \coloneqq \infty$. Then

\begin{enumerate}
    \item  $\mu^*$ is an outer regular measure
    \item All the sets in $\mathcal{R}$ are $\mu$-measurable.
    \item $\mu^*$ coincides with $\mu$ on $\mathcal{R}$
    \item $\mu$ extends as a measure on $\mathcal{M}_{\mu}$.
\end{enumerate}
\label{thm:caratheodory}
\end{theorem}
\begin{proof}
See \cite[1.11.8~Theorem]{Bo07:MeasureTheory}.
\end{proof}

From now on, let $X$ be a topological space.

\begin{definition}[Borel $\sigma$-algebra, Borel measure]
The \emph{Borel $\sigma$-algebra} is the smallest $\sigma$-algebra containing all open sets in $X$. A \emph{Borel measure} is a measure on $\mathcal{B}(X)$.
\end{definition}

\begin{definition}[Support of a Borel measure]
Let $\mu$ be a Borel measure on $X$. The \emph{support} of $\mu$ is defined as the set of all points $x \in X$ for which every open neighbourhood $N_x$ of $x$ has positive measure. We will denote it $\supp(\mu)$.
\label{def:support}
\end{definition}

\begin{definition}[Atom of a measure]
Given a measure space $(X,\mu)$, and $A\subset X$ a measurable subset, we say $A$ is an atom of $\mu$ if the following two conditions are satisfied:
\begin{enumerate}
    \item $\mu(A) > 0$
    \item For any measurable set $B \subset A$ properly contained, $\mu(B)=0$.
\end{enumerate}
\label{def:atommeasure}
\end{definition}


Let $X$ be a locally compact and thus Hausdorff space. A Borel measure $\mu$ on $X$ is  is \emph{locally finite} if every point has a neighborhood $U$ with $\mu(U)< \infty$. The measure $\mu$ is \emph{inner regular} if one has
\[
\mu(U) = \sup \{ \mu(K) : K \subset U \mbox{ compact } \}
\]
for every $U$ open.

For a Borel measure, being outer regular (see Definition~\ref{def:outermeasure}) can be phrased as
\[
   \mu(B) = \inf \{ \mu(U) : U \supset B \mbox{ open } \}
\]
for every Borel set $B$. The measure $\mu$ is \emph{Radon} if it is locally finite, inner regular and outer regular. The Riesz Representation theorem yields a characterization of Radon measures in terms of the associated functional on $C_c(X)$:

\begin{theorem}
  Let $X$ be a locally compact Hausdorff space. For any positive linear functional $F$ on $C_c(X)$ there is a unique Radon measure $\mu$ on $X$ so that
   \[
     F(\psi) = \int_X \psi(x)\mu(x)
  \]
  for all $\psi \in C_c(X)$, and a functional $F$ is called positive if $F(\psi) \geq 0$ for every function $\psi \in C_c(X) \geq 0$ and $x \in X$. The functional $F$ is not assumed to be continuous.
  \label{thm:Riesz}
\end{theorem}

On a general topological space $X$, not all Borel measures are Radon.
If $X$ is a metric space, the Radon property follows from the property of tightness

\begin{definition}
  A \emph{Borel measure} $\mu$ is \emph{tight} if, for every $\epsilon>0$, there exists a compact set $K_{\epsilon}$ so that $\mu(X \backslash K_{\epsilon}) < \epsilon$
\end{definition}

\begin{proposition}
Let $X$ be a metric space.
A tight Borel measure $\mu$ on $X$ is Radon.
\end{proposition}
\begin{proof}
  See \cite[Proposition~2.1.8]{Bogachev18:WeakConvergence}.
\end{proof}

However, tightness and Borel is not equivalent to Radon property: indeed, using the axiom of choice one can construct a separable metric space with a Borel probability measure which is not Radon (see \cite[Example~1.1.2]{Bogachev18:WeakConvergence}).
On complete metric spaces, there are examples of Borel measures which are not Radon. See \cite[Proposition~4.1.6]{Bogachev18:WeakConvergence}.
However, the following Theorem of Ulam will play an important role.

\begin{theorem}
On a complete separable metric space all Borel measures are Radon.
\label{thm:Ulam}
\end{theorem}
\begin{proof}
  See \cite[Theorem~2.1.9]{Bogachev18:WeakConvergence}.
\end{proof}

\begin{definition}
  By the \emph{space of measures $\mathcal{M}(X)$} on a topological space $X$ is the set of all positive Radon measures on $X$.
\end{definition}

We define the following topology in the space of measures.

\begin{definition} The weak$^*$-topology on $\mathcal{M}(X)$ is the weakest topology with respect to which the map
\[
\mathcal{M}(X) \to \mathbb{R}, \mu \mapsto \int \psi \mu
\]
is continuous for all $\psi \in C_c(X)$.
\end{definition}

We suppose that $\mathcal{M}(X)$ is endowed with the weak$^*$-topology.
This topology might seem scary (see Remark~\ref{rmk:signedcurrents}).
However, if $X$ is locally compact, Theorem \ref{thm:Riesz} implies $\mathcal{M}(X)$ is Hausdorff.
In fact, even more is true.

\begin{theorem}
Let $X$ be locally compact. If $X$ is second countable, then $\mathcal{M}(X)$ is metrizable and second countable.
\label{thm:spacemeasureswell}
\end{theorem}

In fact, $\mathcal{M}(X)$ has the Heine-Borel property.
A set $A \subset \mathcal{X}$ is \emph{bounded} whenever
\[
  \sup_{\mu \in H} \int \psi \mu < \infty
\]
for every $\psi \in C_c(X)$.

\begin{theorem}
Let $X$ be a locally compact space. A closed set $A \subset \mathcal{M}(X)$ is compact if and only if it is bounded.
\label{thm:heineborelproperty}
\end{theorem}

In particular, the set of bounded measures on a locally compact space is compact.

\section{Geodesic currents}

Geodesic currents are measures. There are several competing viewpoints on where one can see this measures supported on.
We'll work with all of them, but let us informally discuss the two most widely used in the literature, to our understanding, before we get to the precise definitions.
One viewpoint is as measures supported on the space of geodesics of a surface.
Another is as measures supported on the unit tangent bundle of the surface invariant under the flow.
Both viewpoints have their cons and pros.
The first is more algebraic, and can be easily generalized to hyperbolic groups.
However, the space of geodesics is not well behaved (not Hausdorff), so one has to lift geodesics to the universal cover and work with $\pi_1(\Sigma)$-invariant measures.
The second is more geometric, and allows to see geodesic currents as measures supported on $UT\Sigma$, which is a complete separable metric space. Thence, we don't need to consider the $\pi_1(\Sigma)$-invariance. On the other hand, though, we need to pay the price of considering measures invariant under the geodesic flow. Also, although this definition can be generalized to hyperbolic groups (see \cite[Section~3]{EPS:CountingCurves}) it takes a bit more work to define what unit tangent bundle is in that case (see Section~\cite{EPS:CountingCurves} for details on this).


\begin{definition}[Geodesic current definition 1]\label{def:currents-1}
We define $\GC^+(S)$, the space of oriented geodesic currents on $S$,
to be the
space of $\pi_1(S)$-invariant (positive) Radon measures on
$G^{+}(S)$.
\end{definition}

Since the action of $\pi_1(S)$ on $G^+(S)$ is not discrete, by \ref{prop:fundminimal}, this
definition is hard to visualize. We give alternate
definitions that play a key role in our proofs. For a hyperbolic
surface~$\Sigma$, let $UT\Sigma$ be the unit tangent bundle and let
$\phi_t$ be the geodesic flow on it.

\begin{definition}[Geodesic current definition 2]\label{def:currents-2}
  We can also define oriented geodesic currents to be the (positive) finite
  Radon measures $\mu$ on the unit tangent bundle
$UT \Sigma$ which are invariant under the geodesic flow, in the sense
that $(\phi_t)_*(\mu) = \mu$ for all $t\in\R$. We note that by Theorem~\ref{thm:Ulam}, it suffices to consider Borel measures, since $UT\Sigma$ is a complete separable metric space if $\Sigma$ is a compact surface.
\end{definition}

A measure as defined in Definition~\ref{def:currents-2} induces a flux
measure on transversals to the geodesic foliation of the unit tangent
bundle. This provides an
alternative definition of geodesic currents, which we will exploit in the sequel.

\begin{definition}[Geodesic current definition 3]\label{def:currents-3}
A geodesic current is a family of measures $\{\mu_{\tau}\}_{\tau}$,
where $\tau \subset UT \Sigma$ is a submanifold of the unit tangent
bundle of real codimension 1 transverse to the geodesic
foliation~$\mathcal{F}$, with the following invariance property:  if
$x_1 \in \tau_1, x_2 \in \tau_2$ are two points on transversal
submanifolds on the same leaf of $\mathcal{F}$, and $\phi \colon U_1
\to U_2$ a holonomy diffeomorphism between neighborhoods of $x_1$ and
$x_2$ respectively, then $\phi_{*}\mu_{\tau_1}=\mu_{\tau_2}$.
\end{definition}


\subsection{Oriented vs unoriented currents}

We will be mostly working on the setting of oriented geodesic
currents, but most of our natural examples use unoriented currents.


\begin{definition}[Unoriented geodesic currents]
\label{def:unorientedcurrents}
Let $\GC(\Sigma) \subset \GC^{+}(\Sigma)$ denote the space of \emph{unoriented geodesics currents} defined by
\[
\GC(\Sigma) = \{ \mu \in  \GC^{+}(\Sigma) \mid \sigma_*(\mu)=\mu \}
\]
where $\sigma_* : \mathcal{G}\mathcal{C}^{+}(\Sigma)
\to\mathcal{G}\mathcal{C}^{+}(\Sigma)$ is the map on currents induced
by the \emph{flip map} $\sigma: G^{+}(\Sigma) \to G^{+}(\Sigma)$
that switches the two factors in the definition of $G^+(\Sigma)$,
reversing the orientation of the geodesic.
\end{definition}

There is a map $\Pi : \GC^{+}(\Sigma) \to\GC^{+}(\Sigma)$ 
given by $\Pi(\mu)\coloneqq
\frac{1}{2}(\mu + \sigma_*(\mu))$ with image the subset of unoriented currents.

%%% Local Variables:
%%% mode: latex
%%% TeX-master: "Intersections"
%%% End:
